% block: shifted-systems-chapter-heading page:175 type:chapter
\clearpage\section{{移位线性方程组的应用}}\label{chap:shifted-linear-systems}
\setcounter{footnote}{0}
% block: shifted-systems-abstract page:175 type:ordinary-paragraph
\noindent\textbf{{摘要}}\quad
{如 Section\ \ref{sec:computational-physics} (计算物理) 和 Section\ \ref{sec:machine-learning} (机器学习) 所见, 需要求解一组下列形式的线性方程组}
\[(A+\sigma_k I)\bm x^{(k)}=\bm b,\quad k=1,2,\ldots,m,\]
{其中 $\sigma_k$ 为标量, 这些方程组称为移位线性方程组. 使用 Section\ \ref{subsec:lu-decomposition} 中的 LU 分解时, 需要进行 $m$ 次 LU 分解. 类似地, 使用 Section\ \ref{sec:stationary-iterations} 中的定常迭代方法时, 需要求解 $m$ 个线性方程组, 即对所有 $k$ 求解 $(L+D+\sigma_k I)\bm z=\bm v$. 为求解这些线性方程组, 可以对每个线性方程组应用适当的 Krylov 子空间方法. 另一方面, 若第 $i$ 个线性方程组的初始残差向量与第 $j$ 个线性方程组的初始残差向量共线, 即两个初始残差向量之间的夹角为 $0$ 或 $\pi$, 则生成的 Krylov 子空间相同, 这称为 Krylov 子空间的移位不变性. 这意味着可以共享仅一个 Krylov 子空间的信息来求解所有移位线性方程组, 从而高效计算移位线性方程组. 本章从 Chapter\ \ref{chap:krylov-classification} 中的 Krylov 子空间方法出发, 推导利用移位不变性的 Krylov 子空间方法.}
